% 第 13 题三棱锥正视图：按固定源 PDF 的矢量路径坐标换算重绘。
% 等腰三角形（底 2√3 分两段、高 1），顶点到底边中点的竖直实线；左侧竖直尺寸线（数字 1 断开）；
% 底边下方两条等长水平尺寸线（各被 √3 断开）；底边三处竖直端线。数字只保留原图印出的 1 与两个 √3。
\begin{tikzpicture}[every node/.style={font=\small},line cap=round,line join=round]
% 三角形与中垂实线
\draw (0,0)--(4.30,0)--(2.15,1.24)--cycle;
\draw (2.15,0)--(2.15,1.24);
% 左侧竖直尺寸线（高地：数字 1 断开，上下外向箭头，两端水平端线）
\draw[<->] (-0.38,0)--(-0.38,1.24);
\draw (-0.50,0)--(-0.26,0); \draw (-0.50,1.24)--(-0.26,1.24);
\node[fill=white,inner sep=1.2pt] at (-0.38,0.62) {\(1\)};
% 底边下方两条水平尺寸线（各被 √3 断开，两端外向箭头）
\draw[<->] (0.03,-0.25)--(2.12,-0.25);
\node[fill=white,inner sep=1.2pt] at (1.075,-0.25) {\(\sqrt{3}\)};
\draw[<->] (2.18,-0.25)--(4.27,-0.25);
\node[fill=white,inner sep=1.2pt] at (3.225,-0.25) {\(\sqrt{3}\)};
% 底边三处竖直端线（底边两端与中垂线足）
\draw (0,-0.07)--(0,-0.37);
\draw (2.15,-0.07)--(2.15,-0.37);
\draw (4.30,-0.07)--(4.30,-0.37);
% 图注
\node[below] at (2.15,-0.60) {正视图};
\end{tikzpicture}
